% !TeX spellcheck = es_ES
\documentclass[titlepage,11pt]{article}

\usepackage[spanish,mexico]{babel}
\input{footmult.tex}
\usepackage{graphicx,tabularx,multirow}

\usepackage[onehalfspacing]{setspace}

% HYPERREF
\usepackage{xcolor}
\definecolor{visigrey}{rgb}{.1,.15,.15}
\usepackage[
	allcolors=visigrey,
	colorlinks=true,
]{hyperref} %Hyperref antes de geometry!

\usepackage[a4paper,left=2cm,right=2cm]{geometry}
\geometry{top=1cm,bottom=.5cm}
\savegeometry{titlepage}
\geometry{top=2cm,bottom=2cm}
\savegeometry{main}


\def\bspace{\(\qquad\qquad\qquad\)}
\usepackage[T1]{fontenc}
\usepackage[utf8]{inputenc}
\usepackage{tgheros}
\renewcommand*\familydefault{\sfdefault}

\usepackage[square,numbers,sort]{natbib} % sort=en orden de aparición
\bibliographystyle{unsrtnat}

\setcounter{tocdepth}{2}
% FOOTNOTE

% GLOSARIO
\usepackage[sort=none,abbreviations]{glossaries-extra}
\setabbreviationstyle{long-short}
\newabbreviation{MEM}
{MEM}
{Modal effective mass. Ver \cite{irvine2013effective}}

\newabbreviation{fem}
{FEM}
{modelo de elementos finitos}

\newglossaryentry{trabajo}
{
	name=\emph{Trabajo},
	description={Las actividades comprendidas por la jornada laboral.}
}

\newglossaryentry{mem}
{
	name=Fracción de masa modal efectiva,
	text=MEMF,
	description={Fracción de masa de un sistema que participa del modo vibracional. Ver \cite{irvine2013effective}.}
}

\def\Matlab{\(\textrm{\textsc{Matlab}}\)} %^\textrm{®}

\def\courtesy{Cortesía de Satellogic.}
\def\autor{Patricio Whittingslow}
\def\tema{Práctica Laboral}
\def\titulo{Diseño y análisis\\ estructural de piezas\\ para uso aeroespacial}
\def\legajo{55423}
\def\carrera{Ingeniería Mecánica}
\def\empresa{Satellogic}
\def\fecha{\today}

\begin{document}

    \input{caratula.tex}

\begin{abstract}

Un resumen debería no incluir referencias y no superar 250 palabras.
\end{abstract}

\newpage
\tableofcontents
\newpage
\printunsrtglossaries % list all entries

\vfill
\hrule
\vspace{-1cm}
\subsection*{Nota del autor}
Si el documento es abierto desde un lector PDF moderno (Chrome, Adobe Acrobat Reader, y más), se podrá navegar el mismo mediante las referencias (haciendo click derecho).
\newpage


\section{El objetivo}

Cuando comenzó el \gls{trabajo} el 22 de mayo 2019 la Empresa estaba realizando un estudio de factibilidad para el programa \gls{mem}.

\section{Trabajo}
 
\subsection[Primera etapa (Departamento investigación)]{Primera etapa - Comienzo de \gls{trabajo}} 
Bla Bla.

\subsection[Segunda etapa (Departamento ingeniería)]{Segunda etapa - Dibujo y diseño de bandejas y otras estructuras}

Bla bla



\section{Conclusión}




\section{Opinión de formación recibida}




\section{Anexos}

\addcontentsline{toc}{subsection}{\listfigurename}
\listoffigures
\addcontentsline{toc}{subsection}{\refname}
\bibliography{biblio}

\end{document}